\documentclass[12pt,a4paper]{article}
\usepackage{graphicx}
\usepackage{amsmath}
\usepackage{amssymb}
\usepackage{amsthm}
\usepackage{bm}
\usepackage{hyperref}
\usepackage{booktabs}
\usepackage{geometry}
\usepackage{enumerate}
\geometry{margin=2.5cm}
\hypersetup{colorlinks=true, linkcolor=blue, citecolor=blue, urlcolor=blue}

\newtheorem{theorem}{Theorem}
\newtheorem{proposition}{Proposition}
\newtheorem{corollary}{Corollary}
\newtheorem{lemma}{Lemma}

\begin{document}

\title{Even-$n$ Torus Knots in CDFD:\\
The Exclusion of $T(2,2)$, the Extension to $T(2,2m)$,\\
and the Anti-Phase Mass Scaling Law}

\author{Steve Bico Mujjabi, MD\\
\small Independent Researcher\\
\small ORCID: \href{https://orcid.org/0009-0001-0556-5516}{0009-0001-0556-5516}}
\date{May 2026}

% ─────────────────────────────────────────────────────────────────────────────
\begin{abstract}
Paper~VIII left four open problems, of which this paper resolves three:
the physical identification of the universal anti-phase modes, the
extension of the CDFD torus knot hierarchy to even~$n$, and the
derivation of a mass scaling law for the lightest mode of each family.

The central structural result is the \textbf{Fourier exclusion theorem}:
the universal mass sum rule $\sum_k m_k = 2nM_n$ rests on the Fourier
identity $\sum_{k=0}^{n-1}\cos^2(\theta + 2\pi k/n) = n/2$, which holds
if and only if $n \geq 3$.
For $n = 2$ the identity fails (the residual is $\cos(2\theta)$, not zero),
so the $T(2,2)$ Hopf fibration is excluded from the CDFD framework by
algebraic necessity.

For even $n \geq 4$, the Fourier identity holds and the entire machinery
of Papers~VI--VIII extends without modification.
Self-consistency solutions $\theta_n$, amplitude splits, mode mass spectra,
and the $n^{7/4}$ hierarchy are computed for $n = 4, 6, 8, 10$.
The amplitude-split formula of Paper~VIII applies to even~$n$ with one
exception at $n = 8$.

The anti-phase modes (the universally light pair present in every family
$n \geq 5$) are identified as \textbf{boundary modes}: their mass is set
by how close the self-consistency orientation $\theta_n$ places the two
anti-phase Fourier modes to the amplitude zeros at
$\cos\phi = -1/\sqrt{2}$.
The minimum amplitude satisfies:
\[
  |A_\text{min}| \;=\; \bigl|1 + \sqrt{2}\cos\!\bigl(\theta_n + 2\pi\lceil 3n/8\rceil/n\bigr)\bigr|\,,
\]
which depends on the fractional part $\{3n/8\}$ and is not a simple
power of $n$ or $\theta_n$.
This explains why the two lightest modes are not Goldstone bosons and
have no simple mass formula independent of the knot topology.

A complete mode-mass census for all $T(2,n)$ families with
$3 \leq n \leq 13$ (odd and even) is assembled and presented as a
unified hierarchy table.

\medskip
\noindent In the extended notation used across the series, the vacuum limit is written
as $\Psi_s = (\Phi/C)\cdot S\cdot M_s$. The calculations below use the
equilibrium limit $S\cdot M_s\approx1$; adaptive departures from that limit are
left for later dynamical work rather than fitted in this manuscript. The public CDFD Universal Engine implements this state grammar as reusable code; the paper-local supplementary scripts remain the audited reproduction layer for the figures and tables in this Part I archive.

\end{abstract}

\maketitle

% ─────────────────────────────────────────────────────────────────────────────
\section{Introduction}

The CDFD series~\cite{MujjabiI2026,MujjabiII2026,MujjabiIII2026,%
MujjabiIV2026,MujjabiV2026,MujjabiVI2026,MujjabiVII2026,MujjabiVIII2026}
has so far treated only odd-$n$ torus knot families $T(2,n)$.
Paper~VIII identified even-$n$ families as an open problem and noted
that the universal anti-phase light modes require interpretation.

The present paper addresses both, plus the status of $T(2,2)$.
The torus knot classification~\cite{Rolfsen1976} and the
Faddeev--Niemi energy hierarchy~\cite{FaddeevNiemi1997} underpin
the entire $T(2,n)$ framework; knot polynomial
invariants~\cite{Jones1985} distinguish families that share the
same crossing number, and numerical knot soliton
spectra~\cite{BattyeSutcliffe1998} provide independent checks
on the energy ordering.
Section~\ref{sec:exclusion} proves that $n = 2$ is excluded from the
CDFD framework by a purely algebraic obstruction: the Fourier identity
that underpins the mass sum rule fails for $n = 2$ and only $n = 2$.
Section~\ref{sec:even_n} extends the full hierarchy to even $n \geq 4$.
Section~\ref{sec:boundary_modes} identifies the anti-phase light modes
as boundary modes and derives their mass formula.
Section~\ref{sec:census} presents the complete $T(2,n)$ mass census for
$3 \leq n \leq 13$.

\subsection{Mujjabi boundary principle}

This paper is the first hard boundary test of the hierarchy. A serious theory
must say no. The Mujjabi Boundary Principle is:
\begin{quote}
A candidate CDFD state is physical only if it survives topology, capacity,
phase, and stability constraints simultaneously.
\end{quote}
The exclusion of $T(2,2)$ is therefore not a footnote. It is proof that the
framework is not allowed to accept every attractive pattern. If later candidate
states require exceptions to this boundary principle, the hierarchy loses
predictive force.

% ─────────────────────────────────────────────────────────────────────────────
\section{The Fourier Exclusion Theorem: Why $n = 2$ is Inadmissible}
\label{sec:exclusion}

\subsection{The Fourier identity and where it fails}

The universal mass sum rule (Paper~VII, Theorem~1) rests on a single
algebraic identity:
\begin{equation}
  \sum_{k=0}^{n-1}\cos^2\!\bigl(\theta + \tfrac{2\pi k}{n}\bigr) \;=\; \frac{n}{2}
  \quad\text{for all }\theta\,,
  \label{eq:fourier_id}
\end{equation}
which gives $\sum_k A_k^2 = 2n$ and hence $\sum_k m_k = 2nM_n$.

\begin{theorem}[Fourier exclusion]
  \label{thm:exclusion}
  Identity~\eqref{eq:fourier_id} holds for all $\theta$ if and only if
  $n \geq 3$.
  For $n = 2$ it fails by a theta-dependent residual $\cos(2\theta)$.
\end{theorem}

\begin{proof}
  Writing $\cos^2 x = \tfrac{1}{2} + \tfrac{1}{2}\cos(2x)$:
  \[
    \sum_{k=0}^{n-1}\cos^2\!\bigl(\theta + \tfrac{2\pi k}{n}\bigr)
    = \frac{n}{2} + \frac{1}{2}\sum_{k=0}^{n-1}\cos\!\bigl(2\theta + \tfrac{4\pi k}{n}\bigr)\,.
  \]
  The second sum is the real part of $e^{2i\theta}\sum_{k=0}^{n-1}e^{4\pi ik/n}$.
  This geometric sum vanishes if and only if $e^{4\pi i/n} \neq 1$,
  i.e.\ $4\pi/n \notin 2\pi\mathbb{Z}$, i.e.\ $n \nmid 2$.
  The only integer $n \geq 2$ that divides~2 is $n = 2$ itself.
  For $n = 2$: $e^{4\pi i/2} = e^{2\pi i} = 1$, so the sum equals
  $\sum_{k=0}^{1}e^{2i\theta} = 2e^{2i\theta}$, giving residual
  $\operatorname{Re}(2e^{2i\theta}) = 2\cos(2\theta)$.
  Hence $\sum_k \cos^2 = 1 + \cos(2\theta) \neq 1$ in general.
\end{proof}

\begin{figure}[ht]
  \centering
  \includegraphics[width=0.8\textwidth]{outputs/paper_IX/fig9_0_panel.pdf}
  \caption{The sum $\sum \cos^2(\theta + 2\pi k/n)$ plotted against orientation $\theta$. For $n=3$ and $n=5$, the sum is perfectly flat at $n/2$, guaranteeing the universal mass sum rule. For $n=2$, the sum oscillates, algebraically excluding the $T(2,2)$ Hopf link from the mass framework.}
  \label{fig:fourier_identity}
\end{figure}

\subsection{Physical consequence: $T(2,2)$ is excluded}

For $n = 2$, the sum $\sum_k A_k^2 = 2 + 2\cos(2\theta)$ is
$\theta$-dependent, ranging from $0$ (at $\theta = \pi/2$) to $4$
(at $\theta = 0$).
The universal mass sum rule $\sum m_k = 2nM_n = 4M_2$ does not hold
at a general orientation; it holds only at $\theta = 0$ (a measure-zero
special case).

The self-consistency equation $2\theta_2 = \widetilde{Q}_2(\theta_2)$
does have a numerical solution near $\theta_2 \approx 22.7^\circ$,
but at this $\theta$, $\sum A_k^2 \approx 5.40 \neq 4 = 2n$.
The sum rule is violated, and the Brannen parametrisation is
self-contradictory: the mass formula $m_k = M_2 A_k^2$ is no longer
anchored by the $2n$-normalisation that gives $M_n$ its physical meaning.

\begin{corollary}
  The Hopf fibration $T(2,2)$ is not a member of the CDFD torus knot
  hierarchy.
  The minimal valid family is $T(2,3)$ (the trefoil, the lepton sector).
\end{corollary}

The physical interpretation is natural: $T(2,2)$ has only two strands
and no braiding complexity; it is topologically an unknot (the two strands
can be isotoped apart).
The CDFD framework requires genuine knotting, which begins at $T(2,3)$.

% ─────────────────────────────────────────────────────────────────────────────
\section{Extension to Even $n \geq 4$}
\label{sec:even_n}

\subsection{Validity}

For $n \geq 4$, the divisibility condition $n \mid 2$ is not satisfied,
so Theorem~\ref{thm:exclusion} confirms the Fourier identity holds and
the entire CDFD machinery extends without modification.

The self-consistency condition $n\theta_n = \widetilde{Q}_n(\theta_n)$
(Paper~VIII, Eq.~(2)) applies verbatim.
Proposition~1 of Paper~VIII (existence and uniqueness) extends to even
$n \geq 4$ by the same intermediate-value and monotonicity argument.

\subsection{Self-consistency solutions for even $n$}

Table~\ref{tab:even_sc} gives the solutions for $n = 4, 6, 8, 10$.

\begin{table}[ht]
  \caption{Self-consistency solutions for even-$n$ torus knot families.
    The $n^{7/4}$ hierarchy and sum rule hold for all entries.}
  \label{tab:even_sc}
  \begin{tabular}{cccccc}
    \toprule
    $n$ & Knot & $\theta_n$ (deg) & $M_n$ (MeV) & $\Sigma_n$ (MeV) & Split \\
    \midrule
    4 & $T(2,4)$ & $4.9366$ & $389.44$ & $3115.51$  & $3+1$ \\
    6 & $T(2,6)$ & $2.4595$ & $527.85$ & $6334.15$  & $5+1$ \\
    8 & $T(2,8)$ & $1.4540$ & $654.96$ & $10479.28$ & $6+2$ \\
   10 & $T(2,10)$& $0.8811$ & $774.27$ & $15485.45$ & $7+3$ \\
    \bottomrule
  \end{tabular}
\end{table}

\subsection{Even-$n$ mode spectra}

\begin{table}[ht]
  \caption{$T(2,4)$ mode masses, $\theta_4 = 4.937^\circ$, $M_4 = 389.44$~MeV.}
  \label{tab:Z4_masses}
  \begin{tabular}{ccc}
    \toprule
    Mode & $m_k$ (MeV) & Sector \\
    \midrule
    $m_0$ & $65.135$   & anti-phase \\
    $m_1$ & $300.418$  & co-phase \\
    $m_2$ & $489.995$  & co-phase \\
    $m_3$ & $2259.960$ & co-phase \\
    \midrule
    Sum & $3115.51$ & $= 8M_4$ \\
    \bottomrule
  \end{tabular}
\end{table}

\begin{table}[ht]
  \caption{$T(2,6)$ mode masses, $\theta_6 = 2.459^\circ$, $M_6 = 527.85$~MeV.}
  \label{tab:Z6_masses}
  \begin{tabular}{ccc}
    \toprule
    Mode & $m_k$ (MeV) & Notes \\
    \midrule
    $m_0$ & $30.655$   & anti-phase \\
    $m_1$ & $63.229$   & anti-phase (near zero) \\
    $m_2$ & $89.995$   & anti-phase (near zero) \\
    $m_3$ & $1443.860$ & co-phase \\
    $m_4$ & $1633.222$ & co-phase \\
    $m_5$ & $3073.194$ & co-phase (heaviest) \\
    \midrule
    Sum & $6334.15$ & $= 12M_6$ \\
    \bottomrule
  \end{tabular}
\end{table}

\begin{table}[ht]
  \caption{$T(2,8)$ mode masses, $\theta_8 = 1.454^\circ$, $M_8 = 654.96$~MeV.
    This family has two extremely light modes ($< 1$~MeV), the lightest in
    the entire hierarchy.}
  \label{tab:Z8_masses}
  \begin{tabular}{cc}
    \toprule
    Mode & $m_k$ (MeV) \\
    \midrule
    $m_0$ & $0.411$ \\
    $m_1$ & $0.432$ \\
    $m_2$ & $112.126$ \\
    $m_3$ & $608.793$ \\
    $m_4$ & $702.804$ \\
    $m_5$ & $2552.933$ \\
    $m_6$ & $2685.863$ \\
    $m_7$ & $3815.918$ \\
    \midrule
    Sum & $10479.28$ \\
    \bottomrule
  \end{tabular}
\end{table}

\begin{table}[ht]
  \caption{$T(2,10)$ mode masses, $\theta_{10} = 0.881^\circ$,
    $M_{10} = 774.27$~MeV.}
  \label{tab:Z10_masses}
  \begin{tabular}{cc}
    \toprule
    Mode & $m_k$ (MeV) \\
    \midrule
    $m_0$ & $13.329$ \\
    $m_1$ & $19.029$ \\
    $m_2$ & $132.737$ \\
    $m_3$ & $227.749$ \\
    $m_4$ & $263.816$ \\
    $m_5$ & $1553.077$ \\
    $m_6$ & $1645.125$ \\
    $m_7$ & $3516.772$ \\
    $m_8$ & $3601.652$ \\
    $m_9$ & $4512.167$ \\
    \midrule
    Sum & $15485.45$ \\
    \bottomrule
  \end{tabular}
\end{table}

\begin{figure}[ht]
  \centering
  \includegraphics[width=0.8\textwidth]{outputs/paper_IX/fig9_1_panel.pdf}
  \caption{Mode mass spectra for the even-$n$ torus knot families $n=4,6,8,10$. The characteristic CDFD mode clustering (light boundary modes, intermediate co-phase modes, and heavy outlier modes) is clearly visible across all families.}
  \label{fig:even_spectra}
\end{figure}

\subsection{Amplitude-split formula for even $n$}

The formula of Paper~VIII, Theorem~2:
$n_- = \lfloor 5n/8\rfloor - \lceil 3n/8\rceil + 1$,
extends to even $n$ with one exception.
Table~\ref{tab:even_split} gives results.

\begin{table}[ht]
  \caption{Amplitude splits for even-$n$ families.
    The formula matches numerical results except at $n = 8$.}
  \label{tab:even_split}
  \begin{tabular}{ccccc}
    \toprule
    $n$ & Formula $n_-$ & Numerical $n_-$ & Match & Note \\
    \midrule
    4  & 1 & 1 & Yes & --- \\
    6  & 1 & 1 & Yes & --- \\
    8  & 3 & 2 & No  & Boundary integer: $\lfloor 5\cdot8/8\rfloor = 5$,
                               $\lceil 3\cdot 8/8\rceil = 3$; mode at boundary is positive \\
    10 & 3 & 3 & Yes & --- \\
    \bottomrule
  \end{tabular}
\end{table}

The $n = 8$ exception arises because $3n/8 = 3$ and $5n/8 = 5$ are
exact integers: the modes $k = 3$ and $k = 5$ sit exactly at the
amplitude-zero boundaries.
At the self-consistency orientation $\theta_8 \approx 1.45^\circ > 0$,
these boundary modes are shifted slightly positive, giving $n_- = 2$
rather than 3.
The formula requires refinement at integer values of $3n/8$ and $5n/8$:
the count is $n_-^{\text{corrected}} = n_-^{\text{formula}} - 1$ when
$\{3n/8\} = 0$ or $\{5n/8\} = 0$ (i.e.\ when $8 \mid n$).

% ─────────────────────────────────────────────────────────────────────────────
\section{The Anti-Phase Modes as Boundary Modes}
\label{sec:boundary_modes}

\subsection{The light modes are not Goldstone bosons}

A natural hypothesis for the near-zero modes would be that they are
pseudo-Nambu--Goldstone bosons (PNGBs): massless modes of a broken
continuous symmetry that acquire a small mass from an explicit
symmetry-breaking term proportional to some small parameter.
If so, their masses should scale as $m \propto M_n \theta_n^2$
(the square of the symmetry-breaking angle).

This is \emph{not} the case.
Numerically, $m_\text{min}/(M_n \theta_n^2)$ varies from $1.9$ ($n=5$)
to $782$ ($n=11$) with no discernible universal law.
The hypothesis is refuted.

\begin{figure}[ht]
  \centering
  \includegraphics[width=0.8\textwidth]{outputs/paper_IX/fig9_3_panel.pdf}
  \caption{The minimum mode mass $m_{\text{min}}$ vs the pseudo-Goldstone scaling expectation $M_n \theta_n^2$. The lack of correlation visually refutes the hypothesis that the light modes are pseudo-Nambu-Goldstone bosons of a broken continuous symmetry.}
  \label{fig:goldstone_refutation}
\end{figure}

\subsection{The boundary mode identification}

The amplitude $A_k = 1 + \sqrt{2}\cos\phi_k$ vanishes at
$\phi_k = 3\pi/4$ and $\phi_k = 5\pi/4$.
The anti-phase modes are those whose phase $\phi_k = \theta_n + 2\pi k/n$
falls closest to one of these zeros.

\begin{theorem}[Boundary mode mass]
  \label{thm:boundary_mass}
  The minimum mode mass in any $T(2,n)$ family ($n \geq 5$) is:
  \begin{equation}
    m_\text{min} \;=\; M_n \cdot A_\text{min}^2\,,
    \qquad
    A_\text{min} \;=\; 1 + \sqrt{2}\cos\!\bigl(\theta_n + 2\pi\lceil 3n/8\rceil/n\bigr)\,,
    \label{eq:boundary_mass}
  \end{equation}
  where $k^* = \lceil 3n/8\rceil$ is the mode index whose phase is
  closest to the lower amplitude zero at $3\pi/4$.
\end{theorem}

The key non-universality is the fractional part $\{3n/8\}$, which cycles
through the values $\{1/8, 3/8, 5/8, 7/8, 1/8, \ldots\}$ as $n$ runs
through odd multiples of 1, and through different values for even $n$.
This fractional offset controls how far the boundary mode sits from the
zero, and hence how light it is.

\subsection{Interpretation}

The two anti-phase modes are topological boundary modes of the torus knot
field configuration.
In the $\mathbb{Z}_n$ Fourier decomposition of the knotted soliton, the
modes at $k^*$ and $n - k^*$ are the ones whose phase most nearly
cancels the background condensate.
They are structurally enforced by the sum rule: because $\sum A_k^2 = 2n$
is fixed and the three (or more) co-phase modes carry large amplitude,
the anti-phase pair must carry correspondingly small amplitude.

This is geometrically analogous to the \emph{nodal modes} of a vibrating
string: a string of length $L$ has a mode at every $2L/n$ wavelength,
and the mode whose wavelength is closest to $2L/3$ (the tripling
resonance) is the one with lowest energy.
The CDFD anti-phase modes play this role in the torus knot field.

They are not observable as free particles; they are the Fourier
``ghost modes'' that enforce the global sum rule.
This resolves Paper~VIII open problem~2.

% ─────────────────────────────────────────────────────────────────────────────
\section{The Complete $T(2,n)$ Mass Census}
\label{sec:census}

Table~\ref{tab:full_census} assembles the complete CDFD torus knot
hierarchy for $3 \leq n \leq 13$, combining the odd-$n$ results of
Papers~VI--VIII with the even-$n$ results of this paper.
All energies derive from the single input $m_e = 0.51100$~MeV, via
$M_3$ (Paper~V) and the Faddeev--Niemi scaling $M_n = M_3(n/3)^{3/4}$
(Paper~IV).

\begin{table}[ht]
  \caption{Complete CDFD torus knot mass census, $3 \leq n \leq 13$.
    $\Sigma_n = 2nM_n$ (exact, Paper~VII).
    $\theta_n$: algebraic for $n=3$; self-consistency fixed point for $n \geq 4$.
    Dagger ($\dagger$): the split formula requires a boundary correction (see Section~\ref{sec:even_n}).}
  \label{tab:full_census}
  \begin{tabular}{ccccccc}
    \toprule
    $n$ & Knot & $\theta_n$ (deg) & $M_n$ (MeV) & $\Sigma_n$ (MeV) & Split & Status \\
    \midrule
    3  & $T(2,3)$  & $12.732$ & $313.86$ & $1883.2$  & $3+0$ & Verified (PDG) \\
    4  & $T(2,4)$  & $4.937$  & $389.44$ & $3115.5$  & $3+1$ & Predicted \\
    5  & $T(2,5)$  & $3.698$  & $460.39$ & $4603.9$  & $3+2$ & Predicted \\
    6  & $T(2,6)$  & $2.459$  & $527.85$ & $6334.2$  & $5+1$ & Predicted \\
    7  & $T(2,7)$  & $1.749$  & $592.54$ & $8295.6$  & $5+2$ & Predicted \\
    8  & $T(2,8)$  & $1.454$  & $654.96$ & $10479.3$ & $6+2^\dagger$ & Predicted \\
    9  & $T(2,9)$  & $1.077$  & $715.44$ & $12878.0$ & $7+2$ & Predicted \\
    10 & $T(2,10)$ & $0.881$  & $774.27$ & $15485.5$ & $7+3$ & Predicted \\
    11 & $T(2,11)$ & $0.742$  & $831.65$ & $18296.2$ & $9+2$ & Predicted \\
    12 & $T(2,12)$ & $0.608$  & $886.62$ & $21279.0$ & $9+3$ & Predicted \\
    13 & $T(2,13)$ & $0.528$  & $942.65$ & $24509.0$ & $9+4$ & Predicted \\
    \bottomrule
  \end{tabular}
\end{table}

\subsection{The $n^{7/4}$ hierarchy: extended}

The scaling law $\Sigma_n = 2M_3 \cdot 3^{-3/4} \cdot n^{7/4}$ (Paper~VII)
holds for all $n \geq 3$, odd and even.
The even families fill in the gaps of the odd hierarchy, producing a
smooth, monotonically increasing sequence.
The dominant pattern is that $\Sigma_n$ grows by roughly
$1100$--$1500$~MeV for each unit step in $n$ across the entire range.

\subsection{Theta convergence: the full sequence}

The supplementary script
computes $\theta_n$ for $3 \leq n \leq 13$.
The convergence $\theta_n \to 0$ is monotone and well-described by the
asymptotic $2\pi^2/[(\pi+4)^2 n^2]$ derived in Paper~VIII, confirming
that the large-$n$ limit applies equally to odd and even families.

\begin{figure}[ht]
  \centering
  \includegraphics[width=0.8\textwidth]{outputs/paper_IX/fig9_2_panel.pdf}
  \caption{Theta convergence for the full sequence $3 \leq n \leq 13$. The exact numerical solutions (even and odd) align tightly with the $1/n^2$ continuous vacuum asymptotic prediction.}
  \label{fig:theta_convergence}
\end{figure}

% ─────────────────────────────────────────────────────────────────────────────
\section{Open Problems Closed and Remaining}
\label{sec:open}

\subsection{Closed}

\begin{enumerate}[(a)]
  \item \textbf{Even-$n$ families (Paper~VIII open problem~4)}:
    resolved by extension of the CDFD framework to all $n \geq 4$.
  \item \textbf{Anti-phase light mode identification (Paper~VIII
    open problem~2)}: identified as boundary modes whose mass is set by
    the fractional offset $\{3n/8\}$, not by any symmetry-breaking
    mechanism.
  \item \textbf{Exclusion of $T(2,2)$}: proved algebraically by
    Theorem~\ref{thm:exclusion}.
\end{enumerate}

\subsection{Remaining}

The single surviving open problem from Papers~VII--IX is:

\begin{enumerate}
\item \textbf{Physical principle selecting the self-consistency
  condition.}
  Papers~VII and VIII established that $n\theta_n = \widetilde{Q}_n(\theta_n)$
  has a unique solution for every $n \geq 3$, with an exact closed-form
  asymptotic.
  For $n = 3$, the condition is motivated by the constancy of $Q_3$.
  For $n \geq 4$, a derivation from the CDFD vacuum equation of state
  is needed.
  This is the final open problem of the series.
\end{enumerate}

% ─────────────────────────────────────────────────────────────────────────────
\section{Conclusion}

The Fourier exclusion theorem (Theorem~\ref{thm:exclusion}) identifies
$n = 2$ as the unique excluded family: the Fourier identity
$\sum \cos^2(\theta + 2\pi k/n) = n/2$, which anchors the entire CDFD
mass framework, fails for $n = 2$ and no other $n$.
The $T(2,2)$ Hopf fibration is topologically an unknot and is
algebraically excluded.

For all $n \geq 4$ (even and odd), the framework is valid.
Complete mode spectra for $n = 4, 6, 8, 10$ are tabulated
(Tables~\ref{tab:Z4_masses}--\ref{tab:Z10_masses}) and the
amplitude-split formula is verified with a single boundary-integer
correction at $n = 8$.

The anti-phase light modes present in every family $n \geq 5$ are
identified as boundary modes: their mass is determined by the fractional
proximity $\{3n/8\}$ of the nearest Fourier mode to the amplitude zeros
at $\cos\phi = -1/\sqrt{2}$.
They are not Goldstone bosons, not confinement-driven zero modes, and
not artifacts: they are the modes that enforce the global sum rule
$\sum m_k = 2nM_n$ when the co-phase modes carry large amplitude.

The complete $T(2,n)$ mass census for $3 \leq n \leq 13$
(Table~\ref{tab:full_census}) constitutes the full CDFD torus knot
prediction table, derived entirely from $m_e$ with zero additional free
parameters.

% ─────────────────────────────────────────────────────────────────────────────
\section*{Acknowledgments}
The author used AI-assisted tools for formatting checks and code-verification
support. The scientific claims, interpretations, and final manuscript decisions
are the author's responsibility.

The author is indebted to the foundational work of
D.~Rolfsen~\cite{Rolfsen1976} (torus knot classification),
V.F.R.~Jones~\cite{Jones1985} (knot polynomial invariants),
R.A.~Battye and P.M.~Sutcliffe~\cite{BattyeSutcliffe1998}
(knot soliton energy spectra),
and L.~Faddeev and A.J.~Niemi~\cite{FaddeevNiemi1997} (knots and particles)
--- whose work underpins the extension to even-$n$ torus knots.

% ─────────────────────────────────────────────────────────────────────────────
\appendix
\section{Supplementary Material}

Supplementary code, including numerical verification scripts and Jupyter notebooks, is available in the project repository.

% ─────────────────────────────────────────────────────────────────────────────
\bibliographystyle{unsrt}
\bibliography{references}

\end{document}
